% ARCHIVED DRAFT: independently audited and incorporated in part_B.tex on 2026-11-19 16:10 UTC.
% The assembled paper is authoritative; this file retains drafting provenance.
\begin{saresult}{B.067}{Proved}{A positive non-Noetherian local inclusion with a nonflat analyser}
There is an explicit proper inclusion $R\subset S$
of zero-dimensional local $\mathbb F_2$-algebras
such that every nilpotent squares to zero,
$R$ is non-Noetherian, and every map
\[
 \Omega(R,S)^{\otimes_R n}\longrightarrow\Omega_n(R,S)
\]
is an isomorphism. Nevertheless $\Omega(R,S)$
is not flat over $R$, and the coefficient ring
of B.064 is the proper ring $S$.
Thus flatness of the analyser or of its coefficient
ring cannot be imposed as a necessary condition
in the remaining non-Noetherian classification.
\end{saresult}
\begin{proof}
Index an infinite forest of full binary trees by
$v=(j,w)$, where $j\ge1$ is a root index and
$w$ is a finite binary word. Write $v0,v1$ for
the children. Define
\[
 S=\mathbb F_2[x_v,y_v:v]/
       (x_v^2,y_v^2,\ y_v-x_vy_{v0}y_{v1}:v).             \tag{B67.1}
\]
This presentation has the following explicit
injective direct-limit description. At level
$n$, take the squarefree commutative polynomial
algebra on the variables $x_v$ of depth less
than $n$ and $y_v$ of depth exactly $n$.
Pass to level $n+1$ by replacing each boundary
variable $y_v$ by the product
$x_vy_{v0}y_{v1}$, retaining the old $x$-variables.
Different boundary variables use disjoint new
variable sets. Hence different squarefree
monomials map to different squarefree monomials,
and the map is injective. Its direct limit is
precisely (B67.1); every relation has been
imposed at a finite level, and every finite
polynomial expression is represented at one.
In particular every $y_v$ is nonzero.

Give $x_v$ degree $-1$ and $y_v$ degree $1$.
All transition maps and relations are homogeneous.
Set
\[
                         R=\bigoplus_{d\le0}S_d.          \tag{B67.2}
\]
This is a subring containing $1$. It is proper,
since each $y_v$ is a nonzero element of degree
$1$. Both rings have residue field $\mathbb F_2$.
Indeed their augmentation ideals consist of
elements whose square is zero: in characteristic
two the square of a finite sum is the sum of
the squares, and every nonconstant squarefree
monomial has square zero. Every element with
constant coefficient one is a unit, with
inverse itself. Thus each augmentation ideal
is the unique prime and maximal ideal.
Write $N$ for that ideal in $R$.

For each $k\ge1$, a product of $k$ nonzero
degree-zero monomials $x_vy_{v0}$ from distinct
tree components is nonzero and belongs to
$N^k$. Hence $N$ is not nilpotent. A finitely
generated ideal generated by square-zero
elements is nilpotent, by the pigeonhole
principle. Therefore $N$ is not finitely
generated, and $R$ is non-Noetherian.

The inclusion is a ring epimorphism. Both
$x_vy_{v0}$ and $x_vy_{v1}$ have degree zero
and belong to $R$, so in $S\otimes_R S$ one has
\[
\begin{aligned}
 y_v\otimes1
 &=y_{v0}\otimes x_vy_{v1}\\
 &=x_vy_{v0}\otimes y_{v1}
 =1\otimes y_v.
\end{aligned}                                             \tag{B67.3}
\]
\]
The same identity is automatic for $x_v\in R$.
Since these elements generate $S$, the two
maps $S\to S\otimes_R S$ agree, proving
$S\otimes_R S\simeq S$.

We need a precise conductor calculation:
\[
                        (R:_S N^k)=R\qquad(k\ge1).        \tag{B67.4}
\]
If $c\notin R$, its positive-degree part is
nonzero and involves only finitely many tree
components. Choose $k$ fresh components and
let $a$ be the product of the degree-zero
monomials $x_vy_{v0}$ in them. Then $a\in N^k$,
and multiplying the positive-degree part of
$c$ by $a$ does not kill it. To check this
last assertion without a domain assumption,
each finite-stage algebra is the tensor
product over $\mathbb F_2$ of its separate
tree-component algebras. A product involving
disjoint components is a pure tensor of
nonzero vectors, hence is nonzero; injectivity
of the transition maps preserves this fact.
Thus $ac$ still has a nonzero positive-degree
part, so is outside $R$. This proves (B67.4).

Put $Q(X)=X^4+X^2=X^2(X-1)^2$. For every
$r=a+n\in R$, with $a\in\mathbb F_2$ and
$n\in N$, one has $n^2=0$, so $Q(r)=0$.
Monic division by $Q$ reduces an arbitrary
$S$-polynomial to degree less than four.
For such a remainder $H$, the identity
\[
                    H(a+n)=H(a)+H'(a)n
\]
shows that it preserves $R$ exactly when
$H(0),H(1),H'(0),H'(1)$ belong to $R$,
using (B67.4) with $k=1$.
Hermite interpolation at the two points $0,1$
is invertible over $\mathbb F_2$: it is the
Chinese remainder isomorphism for the
comaximal ideals $(X^2)$ and $((X-1)^2)$.
Therefore those four conditions say exactly
that all four coefficients of $H$ lie in $R$.
We have proved the split module formula
\[
                  \Omega(R,S)=R[X]_{<4}\oplus Q(X)S[X].  \tag{B67.5}
\]

The multivariable version is equally explicit.
Use in each variable the monic basis
\[
 b_d(X)=X^d\ (d<4),\qquad
 b_d(X)=Q(X)X^{d-4}\ (d\ge4).
\]
A tensor product of these bases is an
$S$-basis of $S[X_1,\ldots,X_n]$.
Every basis term with at least one index
$d_i\ge4$ vanishes on $R^n$, so its coefficient
may be arbitrary in $S$. The remaining
remainder has multidegree less than four in
each variable. For each $a\in\{0,1\}^n$,
take finite differences in any subset $J$
of the variables with increments $n_j\in N$.
Since each $n_j^2=0$, the difference is
\[
                      (\partial_J H)(a)\prod_{j\in J}n_j.
\]
If $H$ preserves $R$, it follows that
$(\partial_JH)(a)N^{|J|}\subset R$: the products of the arbitrary increments $n_j$ generate $N^{|J|}$ as an $R$-ideal.
Equation (B67.4) gives $(\partial_JH)(a)\in R$;
the case $J=\varnothing$ is direct evaluation.
Tensoring the four-coordinate Hermite
isomorphism proves that every coefficient
of this low-degree remainder lies in $R$.
Thus $\Omega_n$ has coefficient module $R$
on the all-low-index basis terms and $S$
on all other basis terms.

The tensor power of (B67.5) has exactly these
same coefficient modules: a product of $R$
summands gives $R$, while a product containing
at least one $S$ summand gives $S$, by (B67.3).
The polynomial tensor map identifies these
summands coefficient by coefficient. Hence
every $\theta_n$ is an isomorphism.
For additional explicitness, the two identities
\[
\begin{aligned}
 Q(X+Y)&=Q(X)+Q(Y),\\
 Q(XY)&=Q(X)Q(Y)+Q(X)Y^2+X^2Q(Y)
\end{aligned}
\]
give the co-operations on $cQ$ for every
$c\in S$; multiplication and ordinary
polynomials give them on all of (B67.5).
Composition and every compatibility identity
follow by the tensor isomorphisms as in B.009.

Finally $S$ is not flat over $R$. If it were,
the local inclusion, whose residue fiber is
nonzero, would be faithfully flat; equivalently
its map of one-point spectra is surjective.
A faithfully flat epimorphism is an isomorphism,
contrary to $y_v\notin R$. The two facts are
the opened Stacks Lemmas 10.39.16 (tag 00HQ)
and 10.107.7 (tag 04VU). Since (B67.5)
contains a split copy of $S$, the analyser
is not flat either. Its coefficient ring is
all of $S$, since $cQ$ has leading coefficient
$c$ for every $c$.

The attached exact program \texttt{check\_B\_forest.py}
checks $1890$ transition monomials and $464100$
monomial products, the complete $256$-relation
tensor presentation of one finite branching level,
Hermite ranks $4,16,64$, and the two displayed
polynomial identities. The finite level has tensor
dimension $20$: the parent-variable difference
vanishes while a leaf-variable difference remains
nonzero. Thus the certificate explicitly does
not claim that a finite truncation is epic.
The infinite assertions follow from the proof above.

The epic tensor shuffle is analogous to the
Lazard example reproduced in Stacks Section
110.36, tag 06RH, which was opened. The
bounded-square branching construction, its
injective direct limit, and all integer-valued
and tensor assertions have been proved here;
\cite{StacksEpicExamples}; no priority claim is made and no omitted
detail of that source is used.
\end{proof}
